\documentclass[11pt]{article}

% Packages for formatting
\usepackage[margin=1in]{geometry} % Standard 1-inch margins
\usepackage{times}                % Times New Roman font (Standard for academic SOPs)
\usepackage[utf8]{inputenc}
\usepackage{setspace}             % For line spacing
\usepackage{parskip}              % Adds space between paragraphs, removes indentation

% Set line spacing to 1.15 for readability without wasting space
\setstretch{1.15}

\begin{document}

% Header Section
\begin{center}
    {\Large \textbf{Statement of Purpose}} \\
    \vspace{1em}
    \textbf{Program:} Master of Science in Computer Science (CS75) \hfill \textbf{Applicant:} Yang Sen (Liam) Lin
\end{center}
\vspace{1em}

\textit{I aspire to bridge the gap between the semantic adaptability of foundation models and the rigorous safety guarantees of control theory, architecting trustworthy autonomous systems through algorithmic innovation and hardware-aware optimization.}

In my recent research accepted to \textbf{CoRL 2025}, titled ``See, Point, Fly,'' I identified a fundamental computational trade-off in robotic autonomy: geometric algorithms provide reliability but lack semantic understanding; conversely, Generative Vision-Language-Action (VLA) models offer adaptability but suffer from hallucinations and high inference latency. My objective at UC San Diego is to explore the interfaces between these paradigms. By leveraging the interdisciplinary ecosystem of the \textbf{Contextual Robotics Institute (CRI)}, I aim to investigate how the provable safety guarantees of \textbf{Prof. Sicun Gao} can constrain the generalist policies of \textbf{Prof. Xiaolong Wang}, contributing to the development of ``Trustworthy Embodied Systems.''

My journey toward this objective began with a focus on \textbf{algorithmic decoupling}. In my CoRL 2025 work, I addressed the depth ambiguity problem in monocular navigation not by adding sensors, but by formulating a \textbf{geometric projection algorithm} that grounds 2D VLM prompts into 3D metric controls. By framing action prediction as a spatial grounding task, I achieved state-of-the-art success rates (93.9\%) in unstructured environments. However, I observed that these ``learning-free'' geometric priors became brittle when facing dynamic, non-linear obstacles. This limitation revealed a critical ``reactivity bottleneck'' caused by the inference latency of large models. To address this, I realized that static heuristics must evolve into \textbf{Generalist Learning Policies} supported by rigorous optimization frameworks—a transition that requires the advanced theoretical curriculum offered by UCSD CSE.

Simultaneously, I recognize that algorithmic innovation cannot exist in a vacuum; it must be deployable. My internship at \textbf{Wolley Inc.} grounded me in \textbf{Systems Programming} and architecture-aware optimization. There, I engineered C-based diagnostic tools to profile CXL memory bandwidth, gaining the low-level intuition needed to optimize I/O efficiency for data-intensive AI workloads. I plan to extend this \textbf{Full-Stack Systems} perspective during my upcoming internship at \textbf{OMRON SINIC X}, where I will work on distilling Robotics Foundation Models for edge deployment. This experience confirmed that enabling physical AI requires not just better models, but a holistic approach spanning from memory protocols to high-level reasoning—making the \textbf{Computer Science} major, with its strength in both Systems and AI, the ideal academic home for my research.

My academic background in \textbf{Mechanical Engineering and Computer Science} provides me with a unique advantage: I approach CS problems with physical intuition. I have actively applied advanced mathematics to engineering challenges, utilizing \textbf{Projective Geometry} to derive non-linear unprojection laws in my CoRL paper. This practical application has prepared me for the theoretical rigor of UCSD’s graduate curriculum.

Specifically, the CS75 curriculum aligns perfectly with my technical roadmap to address the ``Safety-Generalization'' tension. I plan to take \textbf{CSE 276F (Machine Learning for Robotics)} with \textbf{Prof. Hao Su} to master Sim2Real transfer techniques for embodied agents. Concurrently, to strengthen my theoretical foundation, I intend to enroll in \textbf{CSE 257 (Search and Optimization)} with \textbf{Prof. Sicun Gao}, acquiring the tools in non-convex optimization needed to design rigorous safety filters. Furthermore, to tackle the computational bottlenecks of VLM deployment, I am eager to take \textbf{CSE 291P (LLM System Optimization)}, leveraging my system profiling background to realize real-time, efficient physical AI.

I specifically seek to contribute to the \textbf{Contextual Robotics Institute (CRI)}, utilizing its collaborative structure to bridge the gap between CSE and ECE methodologies. I aim to work with \textbf{Prof. Sicun Gao (CSE)} to explore neuro-symbolic control frameworks that serve as ``safety filters'' for probabilistic models. Simultaneously, I hope to collaborate with \textbf{Prof. Xiaolong Wang} to apply these constraints to generalist policies learned from large-scale video data. I believe my background—combining the geometric rigor from my CoRL research with the system optimization skills from my industry experience—will allow me to make meaningful contributions to these labs from day one.

Beyond the laboratory, I carry an ``optimization mindset'' from code to community. As the \textbf{GDSC Chapter Lead}, I directed 24 members to bridge the gap between academic theory and industry practice, organizing workshops that impacted over 40 student developers. I aim to bring this leadership to the UCSD community, fostering collaboration between theoretical CS and applied robotics cohorts. My long-term goal is to lead the R\&D of Ubiquitous Physical AI, transforming from a researcher who builds working prototypes into a \textbf{Computational Systems Architect} capable of defining the safety guarantees required to deploy autonomous agents at a global scale.

\end{document}
